\documentclass[10pt,a4paper]{amsart}
\usepackage[margin=19mm]{geometry}
\usepackage{amsmath,amssymb,mathtools}
\usepackage[scr=rsfs,cal=euler]{mathalfa}
\usepackage{xfrac}
\usepackage[unicode,hidelinks,bookmarks=false]{hyperref}
\newcommand{\ie}{i.e.\ }
\newcommand{\cf}{cf.\ }
\newcommand{\resp}{respectively\ }
\newcommand{\Subsection}{\subsection}
\providecommand{\nobreakdash}{\nobreak\hbox{-}}
\setlength{\emergencystretch}{2em}
\multlinegap0pt
\usepackage{amssymb}
\newtheorem{corollary}{Corollary}
\title{The dib-chromatic number of digraphs}
\author{Nahid Javier-Nol, Christian Rubio-Montiel, Ingrid Torres-Ramos}
\date{Source: arXiv:2411.14248v2 (CC BY 4.0)}
\begin{document}
\maketitle

\noindent\emph{Source and licence.} The dib-chromatic number of digraphs. Version arXiv:2411.14248v2 (CC BY 4.0), distributed by arXiv under the Creative Commons Attribution 4.0 International licence, http://creativecommons.org/licenses/by/4.0/, as recorded in the arXiv licence metadata for this exact version and shown on the abstract page https://arxiv.org/abs/2411.14248v2. Full licence notice: \texttt{licenses/arxiv-2411.14248-CC-BY-4.0.txt}.\par\smallskip

\noindent\textit{Editorial scope.} Counted unit: the article's corollary cor9 (source0.tex lines 202--204). The article's complete original proof of it is delivered with it (source0.tex lines 205--213, one \texttt{proof} environment of 9 lines). No mathematics is written, repaired or re-derived: the statement and the proof are the article's own lines, in the article's own notation, and no retained line is substituted or re-flowed.  Retained uncounted as the source-local context the proof needs: lines 66--68.  Nothing was omitted from any retained span, so the package carries no omission record.  No retained line contains a citation, so the wrapper carries no bibliography and no bibliography entry.  No retained line contains a figure environment, an image-inclusion command or drawing code, so no image is delivered and the package declares no asset dependency.  Editorial formatting only, in addition: an AMS \texttt{amsart} preamble that restates the article's own declarations and macros used by the retained lines, namely the environments \texttt{corollary} on separate counters, together with the standard packages those lines need.  The retained environments print in this document as Corollary 1 in order of appearance, whereas the article prints the same items with the numbering of its own full document.  The complete native source of the article as distributed in the arXiv e-print for this exact version is bundled unchanged as \texttt{sources/168892/source0.tex}.
\begin{equation}\label{eq2}
    dib(D) \leq \Delta+1.
\end{equation}

\begin{corollary}\label{cor9}
Let $\overrightarrow{C}_{2m+1}(J)$ be a circulant tournament of order $2m+1$. If $J=\{1,2,\dots,m\}$, then \[dib(\overrightarrow{C}_{2m+1}(J))=m+1.\]
\end{corollary}

\begin{proof}
The coloring $\varsigma(0)=0$ and $\varsigma(i)=\varsigma(i+m)=i$ for $i\in \{1,2,\ldots,m\}$ defines a $b$-coloring of $\overrightarrow{C}_{2m+1}(J)$ because the vertex $0$ is incident to the vertices  $1,2,\ldots, m$ with colors $1,2,\ldots, m$ respectively, and the vertex $0$ is incident from the vertices  $m+1,m+2,\ldots, 2m$ with colors $1,2,\ldots, m$ respectively. Moreover, the vertex $m+i$ (colored $i$) is incident to the vertices $V_i:=\{m+1+i,m+2+i, \ldots, 2m+i=i-1\}$ with colors $i+1,i+2,\ldots, i-1$. Finally, the vertex $i$ (colored $i$) is incident from the set of vertices $V_i$.
Since $\Delta(\overrightarrow{C}_{2m+1}(J))=m$, by Equation \ref{eq2} the result follows.

%$j+i+\left\lfloor \frac{2m+1}{2}\right\rfloor$ is incident to $i$ and incident from $i+\left\lfloor \frac{2m+1}{2}\right\rfloor$ for each $j\in [\left\lfloor \frac{2m+1}{2}\right\rfloor]$.

%$\overrightarrow{C}_{2m+1}(J)$ because the vertex $j+i+\left\lfloor \frac{2m+1}{2}\right\rfloor$ is incident to $i$ and incident from $i+\left\lfloor \frac{2m+1}{2}\right\rfloor$ for each $j\in [\left\lfloor \frac{2m+1}{2}\right\rfloor]$.

\end{proof}

\end{document}
